\documentclass[10pt,a4paper]{amsart}
\usepackage[margin=19mm]{geometry}
\usepackage{amsmath,amssymb,mathtools}
\usepackage[scr=rsfs,cal=euler]{mathalfa}
\usepackage{xfrac}
\usepackage[unicode,hidelinks,bookmarks=false]{hyperref}
\newcommand{\ie}{i.e.\ }
\newcommand{\cf}{cf.\ }
\newcommand{\resp}{respectively\ }
\newcommand{\Subsection}{\subsection}
\providecommand{\nobreakdash}{\nobreak\hbox{-}}
\setlength{\emergencystretch}{2em}
\multlinegap0pt
\usepackage{amssymb}
\usepackage{mathtools}
\usepackage{float}
\newtheorem{corollary}{Corollary}
\newcommand{\E}{\mathbb{E}}
\newcommand{\KL}{\mathrm{KL}}
\newcommand{\abs}[1]{\lvert #1 \rvert}
\newcommand{\piH}{\pi_H}
\newcommand{\tcrit}{t_{\mathrm{crit}}}
\title{Horseshoe Priors and MDP}
\author{Nick Polson, Vadim Sokolov, Daniel Zantedeschi}
\date{Source: arXiv:2604.01266v1 (CC BY 4.0)}
\begin{document}
\maketitle

\noindent\emph{Source and licence.} Horseshoe Priors and MDP. Version arXiv:2604.01266v1 (CC BY 4.0), distributed by arXiv under the Creative Commons Attribution 4.0 International licence, http://creativecommons.org/licenses/by/4.0/, as recorded in the arXiv licence metadata for this exact version and shown on the abstract page https://arxiv.org/abs/2604.01266v1. Full licence notice: \texttt{licenses/arxiv-2604.01266-CC-BY-4.0.txt}.\par\smallskip

\noindent\textit{Editorial scope.} Counted unit: the article's corollary cor:redundancy (source0.tex lines 371--377). The article's complete original proof of it is delivered with it (source0.tex lines 379--386, one \texttt{proof} environment of 8 lines). No mathematics is written, repaired or re-derived: the statement and the proof are the article's own lines, in the article's own notation, and no retained line is substituted or re-flowed.  No further source-local context is needed: every cross-reference of every retained line resolves inside the two delivered spans.  Nothing was omitted from any retained span, so the package carries no omission record.  No retained line contains a citation, so the wrapper carries no bibliography and no bibliography entry.  No retained line contains a figure environment, an image-inclusion command or drawing code, so no image is delivered and the package declares no asset dependency.  Editorial formatting only, in addition: an AMS \texttt{amsart} preamble that restates the article's own declarations and macros used by the retained lines, namely the environments \texttt{corollary} on separate counters, together with the standard packages those lines need.  The retained environments print in this document as Corollary 1 in order of appearance, whereas the article prints the same items with the numbering of its own full document.  The complete native source of the article as distributed in the arXiv e-print for this exact version is bundled unchanged as \texttt{sources/197649/source0.tex}.
\begin{corollary}[Horseshoe redundancy in the sparse normal means model]\label{cor:redundancy}
For the horseshoe prior in the sparse normal means model with $p_0$ signals of size $\abs{\theta_0} \asymp \tcrit$, the per-observation Bayes redundancy (excess KL risk over the oracle who knows which coordinates are signal) satisfies:
\begin{equation}\label{eq:redundancy}
  R_n = \frac{p_0}{2n}\log n + o\!\left(\frac{p_0 \log n}{n}\right).
\end{equation}
The $o(\cdot)$ term absorbs contributions from null coordinates (super-efficient, contributing $O(\tau^4)$ each) and from the self-information correction at the boundary.
\end{corollary}

\begin{proof}
The oracle who knows the signal set achieves per-observation KL risk $p_0/(2n)$ (the parametric rate for $p_0$ parameters). The horseshoe's cumulative redundancy over the oracle is, by Clarke--Barron:
\[
  \sum_{i=1}^n \Bigl(\E[\KL(P_{\theta_0} \| \hat{P}_i^{\mathrm{HS}})] - \E[\KL(P_{\theta_0} \| \hat{P}_i^{\mathrm{oracle}})]\Bigr)
  = \sum_{j \in \text{signal}} \bigl(-\log\piH(\theta_{0j})\bigr) + \sum_{j \in \text{null}} \bigl(-\log\piH(0)\bigr) + O(1).
\]
The null-coordinate sum is $-\infty$ (each term is $-\infty$), but this merely reflects that the horseshoe \emph{outperforms} the oracle on null coordinates. The signal-coordinate sum is $\sum_{j \in \text{signal}} \log\abs{\theta_{0j}}^2 \asymp p_0 \cdot \log\log n$ at the MDP boundary. Dividing by~$n$ gives the per-observation redundancy~\eqref{eq:redundancy}.
\end{proof}

\end{document}
